\documentclass[11pt,a5paper,landscape]{article}

\usepackage{main}
\usepackage{exercices}
\usepackage{tcolorbox}

\renewcommand{\typedocument}{Exercices}
\renewcommand{\classe}{3e}

\begin{document}

\begin{tcolorbox}[
  title=\textbf{Données},
  colback=headergrey!30,
  colframe=black,
  fonttitle=\normalsize,
  boxrule=0.8pt,
  left=6pt, right=6pt, top=4pt, bottom=4pt
]
\begin{multicols}{2}
\begin{itemize}
  \item Intensité de la pesanteur sur Terre : $g = \SI{9,8}{N/kg}$
  \item Intensité de la pesanteur sur la Lune : $g_{\text{L}} = \SI{1,6}{N/kg}$
  \item Intensité de la pesanteur sur Mars : $g_{\text{M}} = \SI{3,7}{N/kg}$
  \item Intensité de la pesanteur sur Jupiter : $g_{\text{J}} = \SI{25}{N/kg}$
\end{itemize}
\columnbreak
\begin{itemize}
  \item Masse de la Terre : $M_T = \SI{5,97e24}{kg}$
  \item Rayon de la Terre : $R_T = \SI{6,37e6}{m}$
  \item Formule de la force de gravitation : $F = G \times \frac{m_1 \times m_2}{d^2}$
  \item Constante de gravitation : $G = \SI{6,67e-11}{} USI$
\end{itemize}
\end{multicols}
\end{tcolorbox}

\vspace{0.4em}

\begin{multicols*}{2}

\exercice{Poids et masse}
\question{Quelle est le poids d'un astronaute de masse \SI{80}{kg} sur la Lune ? Et sur Jupiter ?}
\question{Un véhicule a une masse de \SI{1000}{kg} et un poids de \SI{3700}{N}. Sur quelle planète se trouve-t-il?}
\question{Une balle de golf a une masse de \SI{0,04}{kg} sur Terre.}
\begin{enumerate}[label=\alph*)]
    \item Quelle est sa masse sur la Lune ? 
    \item Quel est son poids sur la Lune ?
    \item Quels appareils peuvent être utilisés pour mesurer sa masse et son poids sur la Lune ?
\end{enumerate}


\exercice{Diagramme objets-interactions}

Une feuille est coincée contre un mur par une main.

Dessiner le diagramme objets-interactions de ce système en indiquant les interactions qui s'exercent sur la feuille.

\exercice{Voiture}

Une voiture de masse \SI{1200}{kg} est à l'arrêt à l'horizontale.

\begin{enumerate}[label=\alph*)]
    \item Schématiser la situation avec un dessin.
    \item Quelles sont les forces qui s'exercent sur la voiture ?
    \item Calculer le poids de la voiture.
    \item Représenter les forces sur le dessin. (1cm = \SI{2000}{N})
\end{enumerate}

\exercice{Gravitation universelle}

\begin{enumerate}[label=\alph*)]
    \item Rappeler la formule du poids d'un objet en fonction de $m$ et $g$.
    \item A la surface de la Terre, le poids d'un objet est égal à la force de gravitation. Ecrire l'égalité.
    \item En déduire la formule de l'intensité de pesanteur $g$ à la surface de la Terre en fonction de sa masse $M_T$ et de son rayon $R_T$.
    \item Calculer l'intensité de pesanteur à la surface de la Terre en utilisant les données fournies.
\end{enumerate}

\end{multicols*}
\end{document}